\section{Об особенностях реализации}
\label{sec:Conclusion_ImplFeatures}
\tikzsetfigurename{Conclusion_ImplFeatures_}

\subsection{Матричные алгоритмы}

Алгоритмы, описанные выше, удобны с точки зрения реализации тем, что могут быть эффективно реализованы с использованием высокопроизводительных библиотек линейной алгебры, которые эксплуатируют возможности параллельных вычислений на современных CPU и  GPGPU~\cite{Mishin:2019:ECP:3327964.3328503}.
Это позволяет с минимальными затратами получить эффективную параллельную реализацию алгоритма для решения задачи КС достижимости в графах.
Благодаря этому, хотя асимптотически приведенные алгоритмы имеют большую сложность чем, скажем, алгоритм Хеллингса, в результате эффективного распараллеливания на практике они работают быстрее однопоточных алгоритмов с лучшей сложностью.

Далее рассмотрим некоторые детали, упрощающие реализацию с использованием современных библиотек и аппаратного обеспечения.

Так как множество нетерминалов и правил конечно, то мы можем свести представленный выше алгоритм к булевым матрицам: для каждого нетерминала заведём матрицу, такую что в ячейке стоит 1 тогда и только тогда, когда в исходной матрице в соответствующей ячейке содержится этот нетерминал.
Тогда перемножение пары таких матриц, соответствующих нетерминалам $A$ и $B$, соответствует построению путей, выводимых из нетерминалов, для которых есть правила с правой частью вида $A B$.

\begin{example}
Представим в виде набора булевых матриц следующую матрицу:
\[
T_0 = \begin{pmatrix}
\varnothing & \{A\}       & \varnothing & \{B\}       \\
\varnothing & \varnothing & \{A\}       & \varnothing \\
\{A\}       & \varnothing & \varnothing & \varnothing \\
\{B\}       & \varnothing & \varnothing & \varnothing \\
\end{pmatrix}
\]

Тогда:
\begin{alignat*}{7}
& &&T_{0\_A} &&= \begin{pmatrix}
0 & 1       & 0 & 0       \\
0 & 0 & 1       & 0 \\
1  & 0 & 0 & 0       \\
0       & 0 & 0 & 0 \\
\end{pmatrix} \ \ \ \ &&T_{0\_B} &&= \begin{pmatrix}
0 & 0       & 0 & 1       \\
0       & 0 & 0       & 0 \\
0  & 0 & 0 & 0       \\
1       & 0 & 0 & 0 \\
\end{pmatrix}
\end{alignat*}
Тогда при наличии правила $S \to A B$ в грамматике получим:
\[
T_{1\_S} =T_{0\_A} \times T_{0\_B} = \begin{pmatrix}
0 & 0       & 0 & 0       \\
0       & 0 & 0       & 0 \\
0  & 0 & 0 & 1       \\
0       & 0 & 0 & 0 \\
\end{pmatrix}
\]
\end{example}

Алгоритм же может быть переформулирован так, как показано в листинге~\ref{lst:cfpq_mtx_bool}.
Такой взгляд на алгоритм позволяет использовать для его реализации существующие высокопроизводительные библиотеки для работы с булевыми матрицами (например M4RI\sidenote{M4RI~--- одна из высокопроизводительных библиотек для работы с логическими матрицами на CPU. Реализует Метод Четырёх Русских. Исходный код библиотеки: \url{https://bitbucket.org/malb/m4ri/src/master/}. Дата посещения: 30.03.2020.}~\cite{DBLP:journals/corr/abs-0811-1714}) или библиотеки для линейной алгебры (например CUSP~\cite{Cusp}).

\begin{algorithm}
  \floatname{algorithm}{Listing}
% TODO!!!
%\begin{algorithmic}[1]
%\caption{Context-free path querying algorithm. Boolean matrix version}
%\label{lst:cfpq_mtx_bool}
%\Function{evalCFPQ}{$D=(V,E), G=(N,\Sigma,P)$}
%    \State{$n \gets$ |V|}
%    \State{$T \gets \{T^{A_i} \mid A_i \in N, T^{A_i}$ is a matrix $n \times n$, $T^{A_i}_{k,l} \gets$ \texttt{false}\} }
%    \ForAll{$(i,x,j) \in E$, $A_k \mid A_k \to x \in P$}
%        %\Comment{Matrices initialization}
%        %\For{$A_k \mid A_k \to x \in P$}
%          {$T^{A_k}_{i,j} \gets \texttt{true}$}
%        %\EndFor
%    \EndFor
%    \For{$A_k \mid A_k \to \varepsilon \in P$}
%       {$T^{A_k}_{i,i} \gets \texttt{true}$}
%    \EndFor

%    \While{any matrix in $T$ is changing}
%        %\Comment{Transitive closure calculation}
%        \For{$A_i \to A_j A_k \in P$}
%          { $T^{A_i} \gets T^{A_i} + (T^{A_j} \times T^{A_k})$ }
%        \EndFor
%    \EndWhile
%\State \Return $T$
%\EndFunction
%\end{algorithmic}
\end{algorithm}

С другой стороны, для запросов, выразимых в терминах грамматик с небольшим количеством нетерминалов, практически может быть выгодно представлять множества нетерминалов в ячейке матрицы в виде битового вектора следующим образом.
Нумеруем все нетерминалы с нуля, в векторе стоит 1 на позиции $i$, если в множестве есть нетерминал с номером $i$.
Таким образом, в каждой ячейке хранится битовый вектор длины $|N|$.
Тогда операция умножения определяется следующим образом:
\[v_1 \times v_2 = \{v \mid \exists (v,v_3) \in P, \textit{append}(v_1, v_2) \& v_3 = v_3\},\] где $\&$~--- побитовое \texttt{``и''}.

Правила надо кодировать соответственно: продукция это пара, где первый элемент~--- битовый вектор длины $|N|$ с единственной единицей в позиции, соответствующей нетерминалу в правой части, а второй элемент~--- вектор длины $2|N|$, с двумя единицами кодирующими первый и второй нетерминалы.

\begin{example}
Пусть $N = \{S, A, B\}$ и в грамматике есть продукция $S \to A B$. Тогда занумеруем нетерминалы $ (S, 0), (A, 1), (B, 2)$. Продукция примет вид $[1, 0, 0] \to [0, 1, 0, 0, 0, 1]$. Матрица $T_0$ примет вид (здесь <<$.$>> означает, что в ячейке стоит $[0,0,0]$):
\[
T_0 = \begin{pmatrix}
. & [0,1,0]       & . & [0,0,1]       \\
. & . & [0,1,0]       & . \\
[0,1,0]       & . & . & . \\
[0,0,1]      & . & . & . \\
\end{pmatrix}
\]

После выполнения умножения получим:
\[
T_1 = T_0 + T_0 \times T_0 =
\begin{pmatrix}
. & [0,1,0]       & . & [0,0,1]       \\
. & . & [0,1,0]       & . \\
[0,1,0]       & . & . & \cellcolor{lightgray}[1,0,0] \\
[0,0,1]      & . & . & . \\
\end{pmatrix}
\]
\end{example}


На практике в роли векторов могут выступать беззнаковые целые числа.
Например, 32 бита под ячейки в матрице и 64 бита под правила (или 8 и 16, если позволяет количество нетерминалов в грамматике, или 16 и 32).
Тогда умножение выражается через битовые операции и сравнение, что довольно эффективно с точки зрения вычислений.

Для небольших запросов такой подход к реализации может оказаться быстрее: в данном случае скорость зависит от деталей.
Минус подхода в том, что нет возможности использовать готовые библиотеки линейной алгебры без предварительной модификации.
Только если они не являются параметризуемыми и не позволяют задать собственный тип и собственную операцию умножения и сложения (иными словами, собственное полукольцо).
Такую возможность предусматривает, например, стандарт GraphBLAS\sidenote{GraphBLAS~--- открытый стандарт, описывающий набор примитивов и операций, необходимый для реализации графовых алгоритмов в терминах линейной алгебры. Web-страница проекта: \url{https://github.com/gunrock/graphblast}. Дата доступа: 30.03.2020.} и, соответственно, его реализации, такие как SuiteSparse\sidenote{SuiteSparse~--- это специализированное программное обеспечения для работы с разреженными матрицами, которое включает в себя реализацию GraphBLAS API. Web-страница проекта: \url{http://faculty.cse.tamu.edu/davis/suitesparse.html}. Дата доступа: 30.03.2020.}~\cite{Davis2018Algorithm9S}.

Также стоит заметить, что при работе с реальными графами матрицы, как правило, оказываются разреженными, а значит необходимо использовать соответствующие представления матриц (CRS, покоординатное, Quad Tree~\cite{quadtree}) и библиотеки, работающие с таким представлениями.

Однако даже при использовании разреженных матриц, могут возникнуть проблемы с размером реальных данных и объёмом памяти.
Например, для вычислений на GPGPU лучше всего, когда все нужные для вычисления матрицы помещаются на одну карту.
Тогда можно свести обмен данными между хостом и графическим сопроцессором к минимуму.
Если не помещаются все, то нужно, чтобы помещалась хотя бы тройка непосредственно обрабатываемых матриц (два операнда и результат).
В самом тяжёлом случае в памяти не удаётся разместить даже операнды целиком и тогда приходится прибегать к поблочному умножению матриц.

Отдельной инженерной проблемой является масштабирование алгоритмов на несколько вычислительных узлов, как на несколько CPU, так и на несколько GPGPU.

Важным свойством рассмотренного алгоритма является то, что описанные проблемы с объёмом памяти и масштабированием могут эффективно решаться в рамках библиотек линейной алгебры общего назначения, что избавляет от необходимости создавать специализированные решения для конкретных задач.

\subsection{Алгоритм на основе тензорного произведения}

Как и алгоритмы, представленные в разделе~\ref{chpt:MatrixBasedAlgos}, представленный здесь алгоритм оперирует разреженными матрицами, поэтому, к нему применимы все те же соображения, что и к алгоритмам, основанным на произведении матриц. Более того, так как результат тензорного произведения является блочной матрицей, то могут оказаться полезными различные форматы для хранения блочно-разреженных матриц. Вместе с этим, в некоторых случаях матрицу смежности рекурсивного автомата удобнее представлять в классическом, плотном, виде, так как для некоторых запросов её размер мал и накладные расходы на представление в разреженном формате и работе с ним будут больше, чем выигрыш от его использования.


Также заметим, что блочная структура матриц даёт хорошую основу для распределённого умножения матриц при построении транзитивного замыкания.

Вместо того, чтобы перезаписывать каждый раз матрицу смежности входного графа $M_2$ можно вычислять только разницу с предыдущим шагом.
Для этого, правда, потребуется хранить в памяти ещё одну матрицу.
Поэтому нужно проверять, что вычислительно дешевле: поддерживать разницу и потом каждый раз поэлементно складывать две матрицы или каждый раз вычислять полностью произведение.

Заметим, что для решения задачи достижимости нам не нужно накапливать пути вдоль рёбер, как мы это делали в примерах, соответственно, во-первых, можно переопределить тензорное произведение так, чтобы его результатом являлась булева матрица, во-вторых, как следствие первого изменеия, транзитивное замыкание для булевой матрицы можно искать с применением соответствующих оптимизаций.
